% AI-written, AI cross-reviewed; no human sign-off; not a claim that RH is proved
% Manuscript v3, heading p. 9; original out/r23_p017.tex, SHA-256 49f2380c337781b49a96cff47c1372f8d5923ff45d8bf22629de80a824819866.
\subsection{From a tied lattice chain to the level prime sum}\label{r23:p017}

\begin{proposition}[P017, family window estimate]\label{r23:family}
Fix a level $\Lambda$, real $e_D,e_L$ and $\varepsilon>0$. Let $F$ be a
family at $\varepsilon$ with the uniform explicit-formula hypothesis
$\mathrm{EFamC}(\Lambda,F)$. Suppose that there exist $K,c_{HH},e_K,C_L,x_0$
such that, for every $x\ge x_0$ in $\mathrm{Win}(F)$,
$\mathrm{ChainTied}(F;e_D,e_L,K,c_{HH},e_K,C_L,x)$ holds.
Then there exists $X$ such that on this same set for $x\ge X$,
\[
     \mathcal G_\Lambda(x)\le x^{1/2+(e_D+e_L+7)\varepsilon}.
\]
Here $\mathcal G_\Lambda$ is the real weighted prime sum of
Definition~\ref{def:E-level}. The estimate is one-sided.
\end{proposition}
\noindent\nolinkurl{UpperTheory.Comparison₂.level_gm_le_on_family_scales},
\path{LevelGmLeOnFamilyScales.lean:28--30};
\nolinkurl{Formalism.HC.ChebyshevInequalityOnWindows},
\path{TotientGaloisClasses/TopWin.lean:23--28}.

\paragraph{The column scale.}
In the arithmetic construction, $\delta_1=x^{-\varepsilon/8}$. With
$A=\sqrt{x\delta_1}=x^{1/2-\varepsilon/16}$, the construction may take
the even integer $m=2\lceil A/2\rceil$. For $0<\varepsilon<1/52$
and $x\ge16$, $A\ge2$ and $A\le m<A+2\le2A$. This gives the $m$ 
clause of \nolinkurl{GlobalTheory.EX.RankVal} at the prescribed frame, whose
order lies in $[x^{2+8\varepsilon},2x^{2+8\varepsilon}]$.
For $x^{9\varepsilon/16}\ge2$ the resulting bounds imply
\[
    x^{1/2-\varepsilon}\le m\le x^{1/2+\varepsilon/2},\qquad
    x/m\le x^{1/2+\varepsilon}.
\]
Thus dividing a cost of order $x$ by the number of columns produces a
square-root scale. A second cost is proportional to $m$ itself. The four
terms after division, including the density losses, are evaluated in
Lemma~\ref{r23:budget}. In Proposition~\ref{r23:family} the two wider rank
bounds are supplied directly in the tied size windows, so the implication
has the stated quantifier $\varepsilon>0$.
\decid{D513};
\nolinkurl{Scale.delta1}, \path{LocalRingCardinality/DataParams₂.lean:45};
\nolinkurl{GlobalTheory.EX.RankVal}, \path{PairRowIffQdRows.lean:28--32};
\nolinkurl{GlobalTheory.EX.exists_even_between} and
\nolinkurl{GlobalTheory.EX.exists_ranks_rankVal_frame},
\path{ExistsRanksRankValFrame.lean:10--15,45--99};
\nolinkurl{Deduction.ch_core}, \path{LandauTheorems/ChCore.lean:343--351}.

\paragraph{The bound before taking characters.}
At one scale let $H$ be the finite additive group of the reference frame,
$A_0\le H$, $TT\subset H$, $T_2\subset A_0$ and
$n=|H|$, $n_T=|TT|$, $n_2=|T_2|$, $\delta_2=n_2/n$.
Write $\mathrm{gap}(u)$ for the sum of
$\mathrm{ASK}-\Re\mathrm{MAIN}$ over $t\in TT$ and $1\le j\le m$
at translate $u$. Let $B$ be the four-term bracket in
\hyperref[sub:B09]{B09}, with $a_E$ the lattice parameter $a$.
\nolinkurl{Deduction.CH.gapA}, \path{DedekindLFunctions/SplitInertia.lean:50--52};
\nolinkurl{errorBracketSizeWindows}, \path{ErrorFormsWeakWindowsFalse.lean:133--135}.

\begin{formalstatement}[Lattice inputs for the gap]\label{r23:gapinputs}
Under the eleven hypotheses in \nodeid{P021}, at each lattice unit $v$,
$\mathrm{ASK}(v)-\mathrm{Main}(v)\le\mathrm{ErrX5at}(v)$.
Under the error windows and $\mathrm{LargeX}$,
$\mathrm{ErrX5at}(v)\le x^{(e_{\rm Div}+e_{\rm Loss}+3)\varepsilon}
\mathrm{ErrText}$ (\decid{D567}). The actual bounds for $B_0,N_c$
and the required size and identification hypotheses are specified in
\hyperref[sub:B10]{B10}. The weighted mean and slope inputs to P021
are recorded at \decid{D568} and \decid{D570}.
\end{formalstatement}
\noindent\nolinkurl{each_unit_bound_statement}, \path{LevelOfCongruence.lean:175--197};
\nolinkurl{errX5at_le}, \path{ChainSlopeWindows/PChain.lean:613--672}.

\begin{proof}[Derivation of the gap bound, B10]
For $v$ mapping to $u$, replace the top and place indices by the given
bijections with $TT$ and $\Phi$. The summands of the lattice main term
then have exactly the recorded $a_t$, $w_{u,j}$, Euler constant and
$\Re\ell(\chi_C)$, so the sum is $\sum_{t,j}\Re\mathrm{MAIN}_{t,j}$.
The identification of the ask term gives $\sum_{t,j}\mathrm{ASK}_{t,j}$.
Their difference is $\mathrm{gap}(u)$ (\decid{D569}).
Now $\mathrm{ErrText}=n_Ta_E(1+m^2/k)+m(B_0+N_c)$.
Insert $B_0\le2nmbx^\varepsilon$ and
$N_c\le n_2 k b x^\varepsilon/(n^2x^{-\varepsilon})$.
All multipliers are nonnegative, so $\mathrm{ErrText}\le B$.
Combining the two inequalities in \ref{r23:gapinputs} gives
\[
    \mathrm{gap}(u)\le x^{(e_{\rm Div}+e_{\rm Loss}+3)\varepsilon}B.
\]
Every $u\in A_0$ has a preimage by the coverage clause in
\nolinkurl{CoreAtH}; this transfers the bound to the entire subgroup.
\nolinkurl{Deduction.ch_top_main_eq}, \path{LandauTheorems/ChTop.lean:11--35};
\nolinkurl{Deduction.ch_top}, \path{LandauTheorems/ChTop.lean:37--72};
\nolinkurl{CoreAtH}, \path{PrimitiveCosets/CoreAtH.lean:29--41}.
\end{proof}

\paragraph{Which characters are being summed.}
Put $E_0=(\mathbb Z/2)^2$, $s_0=((1,0),0)$,
$e_1(t)=s_0+(0,t)$ and
$\mathcal A=\{\xi\in\widehat{E_0\times H}:\xi(0,A_0)=1\}$.
The character sum at $y$ is denoted by $G_\xi(y)$.
Write $K'$ for the kept subset, $o=\Lambda.o$, $c=(3o+1)^{-1}$, $M=4o$,
$K_o=\sum_{\psi\in\Lambda.\mathrm{KEPT}}|(L'/L)(1,\psi)|$
with the source's convention for the logarithmic derivative at one.
Define $v_\xi=\mathrm{wt}_c(\xi)\Re G_\xi(x)$ for $\xi\in K'$,
and $v_\xi=-|G_\xi(x)|$ for $\xi\in\mathcal A\setminus K'$.
Here $\mathrm{wt}_c$ is $1$ on the characters trivial on $H$, and $c$
on the other characters. Thus
$G:=\sum_{\xi\in\mathcal A}v_\xi=G_m^A(x)-\Sigma_A(x)$
by splitting the finite sum (\decid{D565}).
\nolinkurl{UpperTheory.Comparison₂.koOf}, \path{PrimitiveCosets/CoreAtH.lean:27};
\nolinkurl{Deduction.CH.wtA}, \nolinkurl{GmA}, \nolinkurl{SigA},
\path{DedekindLFunctions/SplitInertia.lean:26--40};
\nolinkurl{CharacterBalanceMoves.sum_vv}, \path{LandauTheorems/ChCore.lean:183--200}.

\begin{formalstatement}[Character inputs at the same host]\label{r23:charinputs}
Let the host satisfy the axioms \nolinkurl{HostAx} on its frame. Assume $\mathrm{EAtH}(\Lambda,H,TT,x,\varepsilon,C_E)$
and $\ell$ agrees with its primitive-character logarithmic derivatives.
Then $0\in K'\subset\mathcal A$, all kept characters take value one at
$s_0$, and $|\mathcal A|\le M$. For kept $\xi$ nontrivial on $H$,
$|\sum_{t\in TT}\xi(0,t)-c n_T|\le M x^{-1+\varepsilon}n_T$.
Nonzero kept characters satisfy $|\ell(\xi)|\le K_o$ and discarded
members of $\mathcal A$ satisfy $|\ell(\xi)|\le C_E\log x$
(\decid{D548}). Also $G_m^A(y)=\mathcal G_\Lambda(y)$ for every $y$
(\decid{D547}) and each discarded $|G_\xi(x)|\le x^{1/2+\varepsilon}$
(\decid{D549}). There is an absolute $C_d\ge0$ such that for $y,z\ge1$
the increments of $G_0$ are bounded by $C_d|y-z|$, and those of each
nontrivial $G_\xi$ by
$|y-z|(\log\max(y,z)+C_d+|\ell(\xi)|)$ (\decid{D562}).
\end{formalstatement}
\noindent\nolinkurl{charsA}, \path{PrimitiveCosets/charsA.lean:80--125};
\nolinkurl{gmA_eq}, \path{PrimitiveCosets/GmAEq.lean:61--149};
\nolinkurl{sigA_le}, \path{PrimitiveCosets/SigALe.lean:38--60};
\nolinkurl{derivative_integral_all}, \path{LipschitzWeightAndMinimum.lean:966--970}.
The identification in D547 sends $\xi$ to the primitive character inducing
its host character. Existence and injectivity give a bijection on the kept
sets. Triviality on $H$ singles out the principal and distinguished quadratic
cases; the exceptional-character field of $\mathrm{EAtH}$ matches their
weights. Compatibility with $\ell$ identifies the prime sums termwise.
For the omitted set the same reading and the host bound give D549 after
summing at most $4o$ terms.

\begin{proof}[Conversion to a character lower bound, D564--D565]
Use $y_t=a_t/b$, so $|y_t-x|\le1$, and take $x\ge3$.
For nonzero $\xi$, $G_\xi(1)=\ell(\xi)$, whence the increment estimate gives
$|G_\xi(x)|\le x(\log x+C_d+|\ell(\xi)|)$.
For nonzero kept $\xi$, first replace $G_\xi(y_t)$ by $G_\xi(x)$ under
$\sum_t\xi(e_1(t))$. The cost per term is at most
$\log(x+1)+C_d+K_o$. Next replace the sum of its coefficients by
$n_T\mathrm{wt}_c(\xi)$: the values at $e_1(t)$ and $(0,t)$ agree,
and the total error is bounded by
$n_T M x^{-1+\varepsilon}x(\log x+C_d+K_o)$.
For $\xi$ trivial on $H$ this latter error is zero. For the principal
frequency the entire cost per term is $C_d$. For each discarded frequency
use $\Re(\xi(e_1(t))G_\xi(y_t))\ge-|G_\xi(y_t)|$ and the increment
bound to get a lower bound $v_\xi-(\log(x+1)+C_d+|C_E|\log x)$.
These are exactly the three cases of \decid{D564}.

For a threshold independent of the host, let $s=x^{1/4}$,
$e=x^\varepsilon$ and
\[
 A_1=5+C_d+|K_o|+|M|(4+C_d+|K_o|)+4|C_E|.
\]
Since $\log x\le4s$, $\log(x+1)\le5s$ and $1\le s,e$, each of the
three per-character costs per $t$ is at most $A_1 s e$ (the coefficient
error uses $x^{-1+\varepsilon}x=e$). Increase the threshold until
$|M|A_1\le s$. Summing over the band costs at most $n_Ts^2e$,
hence also at most $4n_Tx^{1/2+\varepsilon}$. Taking real parts gives
\[
 n_T(G-4x^{1/2+\varepsilon})
 \le \sum_{t\in TT}\Re\sum_{\xi\in\mathcal A}
                 \xi(e_1(t))G_\xi(y_t).
\]
\nolinkurl{CharacterBalanceMoves.move}, \nolinkurl{triv} and \nolinkurl{one_char},
\path{LandauTheorems/ChCore.lean:154--179,202--281};
\nolinkurl{Deduction.ch_core}, \path{LandauTheorems/ChCore.lean:352--359,375--439}.
\end{proof}

\paragraph{Averaging over the translates.}
The formal inputs to the following calculation are the exact host
balance $\mathrm{ASK}-\mathrm{MAIN}=\mathrm{De}$ for positive
$a_t,w$ and $a_t/w\le N$, and the weight increment bound
\[
 |wG_\xi(A/w)-bG_\xi(A/b)|\le
   3A|w-b|/\min(w,b)\quad (A,w,b>0).
\]
These enter the formal statement \decid{D561}. The host assumption for
balance on the excluded set $S=\{q\}\cup\mathrm{small}$ is \decid{D014}.
\nolinkurl{Deduction.ch_balance}, \path{LandauTheorems/ChBalance.lean:200--203};
\nolinkurl{lipschitz_weight_all}, \path{LipschitzWeightAndMinimum.lean:957--964};
\nolinkurl{hostFactsS}, \path{PrimitiveCosets/hostFactsS.lean:127--136}.

\begin{proof}[Finite-group average and B09]
Put $f=|\Phi|=2n$. For each $u\in A_0$, retain the thick places and freeze
their widths at $b$. The resulting Fourier sum for $t\in TT$ is
\[
 T(u,t)=\frac{b}{f}\sum_{s\in T_2}\;
        \sum_{\xi\in\widehat{E_0\times H}} \xi(-e_1(s+u-t))G_\xi(y_t).
\]
To compare with one column of $\mathrm{De}$, write both sums on the
same translated complement of $\Phi$. A thick place there is indexed
once by $s\in T_2$, with the Fourier factor in the display. At a thick
place $b\le w\le b+\eta+1$ and $a_t\le(x+1)b$, so replacing its
weighted value by the frozen one costs at most $3(x+1)(\eta+1)$.
At a thin place $0<w\le2b'$. Since $a_t\ge(x-1)b$,
$b'X'\le xb$ and $x\ge3$, one has $a_t/w\ge X'/3$.
Thinness therefore bounds $w|G_\xi(a_t/w)|$ by $a_t\theta\le
(x+1)b\theta$. There are at most $f$ places and $2f$ characters,
each with unit-modulus Fourier coefficient and normalization $1/f$.
The sum of absolute errors is at most
\[
 2f\big[3(x+1)(\eta+1)+(x+1)b\theta\big]\le 2b/x
\]
by the two cost inequalities in \hyperref[sub:B09]{B09}.
The balance identity now gives
$\Re T(u,t)\le\mathrm{ASK}_{t,j}-\Re\mathrm{MAIN}_{t,j}+2b/x$
for every column $j$. Thus if $\mathrm{gap}(u)\le E$ on $A_0$, then
$\sum_{t,j}\Re T(u,t)\le E+2n_Tmb/x$ on $A_0$.
\nolinkurl{CharacterBalanceMoves.Th}, \nolinkurl{th_le},
\path{LandauTheorems/ChCore.lean:16--20,75--152};
\nolinkurl{CharacterBalanceMoves.core}, \path{LandauTheorems/ChCore.lean:283--320}.

For $s\in T_2\subset A_0$, sum $\xi(-e_1(s+u-t))$ over $u\in A_0$.
Its value is $|A_0|\xi(e_1(t))$ if $\xi\in\mathcal A$, since
$-e_1(s+u-t)=e_1(t)+(0,-s)+(0,-u)$.
Otherwise choose $u'\in A_0$ with $\xi(0,u')\ne1$; translating the
sum by $u'$ multiplies it by $\xi(0,-u')$, hence the sum vanishes.
Consequently the average of $T(u,t)$ over $A_0$ equals
$b(\delta_2/2)\sum_{\xi\in\mathcal A}\xi(e_1(t))G_\xi(y_t)$:
the factor is $bn_2/f$. Sum over all $t$ and the $m$ columns and apply
the preceding character lower bound. Since $|A_0|>0$, the result is
\[
    n_T m b(\delta_2/2)(G-4x^{1/2+\varepsilon})\le E+2n_Tmb/x .
\]
With $E=x^{C\varepsilon}B$, $C\ge0$, the size-window calculation
in Lemma~\ref{r23:budget} now gives $G\le x^{1/2+(C+3)\varepsilon}$.
Its hypotheses are precisely \nolinkurl{RestSize}, the upper rank and
lower $k$ inequalities, $x\ge3$, $26\le x^{3\varepsilon/8}$, and the
sign and cardinality conditions (\decid{D563}, \decid{D566}).
Adding $\Sigma_A$ proves B09. The threshold precedes the choice of
$H,TT,T_2$, the weights, and $m$.
\nolinkurl{CharacterBalanceMoves.sum_A}, \nolinkurl{sum_Th},
\path{LandauTheorems/ChCore.lean:28--73};
\nolinkurl{CharacterBalanceMoves.core}, \path{LandauTheorems/ChCore.lean:322--338};
\nolinkurl{Deduction.ch_core}, \path{LandauTheorems/ChCore.lean:343--370,419--442}.
\end{proof}

\begin{proof}[The one-scale estimate, B08]
Fix the level and all the constants of \hyperref[sub:B08]{B08}. Take the
maximum of three thresholds. First take the B09 threshold with
$C=\max(e_D+e_L,0)+3$, $c=(3o+1)^{-1}$, $M=4o$, $C_E$ the
explicit-formula constant and $K_o$ as above. Here $|c|\le1$ since
$o\ge0$. Second require
$\log x,\,200(1+K)^2(1+c_{HH})(1+e_K)\le x^\varepsilon$
(\decid{D550}). Third require $1+4o\le x^\varepsilon$.
All three thresholds depend only on the fixed constants.

For any host above that threshold with both the explicit-formula and
core packages, write $c'_{HH},e_{\rm Div},e_{\rm Loss}$ for its internal
constants. Its error windows give $K,e_K,c'_{HH},e_{\rm Div},e_{\rm Loss}\ge0$.
Since $c'_{HH}\le c_{HH}$, the same largeness bound holds for $c'_{HH}$.
Also $0\le e_{\rm Div}\le e_D$, $0\le e_{\rm Loss}\le e_L$,
so $C=e_D+e_L+3$. The data have scale $x$ and exponent $\varepsilon$.
The bracket $B$ is nonnegative. Thus B10 and monotonicity of $x^a$
in $a$ give $\mathrm{gap}(u)\le x^{C\varepsilon}B$ for all $u\in A_0$.
The sizes, thinness, costs and containment $T_2\subset A_0$ are fields
of the same core package. D014 supplies its arithmetic host facts and
D548 the required character data. Hence B09 applies.
Using D547 to identify the kept sum and D549 to add the omitted terms,
\begin{align*}
 \mathcal G_\Lambda(x)
   &\le x^{1/2+(e_D+e_L+6)\varepsilon}+4o\,x^{1/2+\varepsilon}\\
   &\le (1+4o)x^{1/2+(e_D+e_L+6)\varepsilon}
    \le x^{1/2+(e_D+e_L+7)\varepsilon}.
\end{align*}
This establishes B08 for all eligible frames and hosts simultaneously.
\nolinkurl{UpperTheory.Comparison₂.core_gm}, \path{PrimitiveCosets/CoreGm.lean:12--74};
\nolinkurl{large_x_eventually}, \path{SubThreeMulLogTwoLe.lean:319--329}.
\end{proof}

\begin{proof}[Family quantifiers and proof of Proposition~\ref{r23:family}]
Unpack $\mathrm{EFamC}$ to obtain $C_E,x_E$ such that every member at
each $x\ge x_E$ satisfies $\mathrm{EAtH}$ with $C_E$. Fix the five
eventual-tiedness witnesses of the proposition. Let $X$ be the maximum
of $x_0,x_E$ and the B08 threshold for these constants. For $x\in
\mathrm{Win}(F)$, $x\ge X$, unpack $\mathrm{ChainTied}$ at this $x$.
Use its reference member $i_0$, take the translation map to be the identity
on its frame group, and use the word fibre of the assigned member
$\mathrm{pt}(u)$ at each $u$. This covers $A_0$. The arithmetic tie's
test-set membership equation and the family's common $TT$ image identify
$TT$ with the reference host's $TT$ by injectivity of the frame coordinate.
Its digit description at translate $0$ expresses each $s\in T_2$ as a digit
in the cyclic subgroup $A_0$. The same tie supplies the compatible
$\ell$. Restrict the weight bounds and cutoff from all translates to
$A_0$, and retain the two bijections, identities, lattice size and error
bounds, thin scales, costs and remaining-size bounds. These are exactly
the existential witnesses for $\mathrm{CoreAtH}$ at $i_0$.
Evaluate the universally quantified explicit-formula bound at this very
$i_0$. B08 now gives the claimed estimate. Since $X$ was fixed before
$x$ and $i_0$, it works on the whole tail of the window.
\nolinkurl{Formalism.HC.EFamC}, \path{TotientGaloisClasses/TopDefs.lean:15--16};
\nolinkurl{ChainTied}, \path{GlobalSchurBounds/TieGlobal.lean:97--114};
\nolinkurl{UpperTheory.Comparison₂.tiedCore_of_chainTied},
\path{LevelGmLeOnFamilyScales.lean:13--26};
\nolinkurl{UpperTheory.Comparison₂.windowInequality_of_artinInputs},
\path{WindowInequalityOfArtinInputs.lean:17--25};
\nolinkurl{UpperTheory.Comparison₂.level_gm_le_on_family_scales},
\path{LevelGmLeOnFamilyScales.lean:28--30}.
\end{proof}
